\section{Proposed Framework}
We design a report-bound zero-knowledge compliance method for multisource spacetime observations subject to measurement uncertainty and Byzantine failures. Instead of first collapsing these observations into a single estimated location, the method retains their intervals and evaluates compliance against every possible state reaching a support threshold. The workflow constructs a robust feasible set from signed evidence, determines its relationship to the authorized policy, and proves the resulting decision. We first establish the system model and decision semantics, and then describe the binding protocol and proof circuits.

\subsection{System model and assumptions}
For a reference event associated with a sensing report, write the true spacetime state as $z^\star=(x,y,t)$. We represent space and time in a $d$-dimensional discrete domain,
\begin{equation}\label{eq:domain}
\Omega=\left(\prod_{a=1}^{d}[L_a,H_a]\right)\cap\mathbb Z^d.
\end{equation}
The interval $[L_a,H_a]$ bounds coordinate $a$. An authorized spacetime policy is a public integer box $P\subseteq\Omega$. Each enrolled evidence issuer $i$ signs an inclusive axis-aligned integer box $U_i\subseteq\Omega$ concerning the same reference event, under a common context $c$. All sources use a common coordinate reference system, time base, measurement units, and quantization rule. In the prototype, interval endpoints are encoded as unsigned 16-bit integers. Exact evaluation consequently refers to the quantized finite domain, not to error-free measurement in continuous physical space.

The system comprises evidence issuers, a vehicle acting as prover, and a verification gateway. The host crowdsensing application supplies the authorized task, raw report bytes, roster of issuers, and calibration settings. The vehicle collects the signed evidence for a report and produces a zero-knowledge proof of its four-valued classification. The gateway verifies the proof under its trusted task configuration and returns a decision. Participant recruitment, incentives, dynamic active-set maintenance, and independent institutional audits of issuers are outside this module.

Assume that among $n$ enrolled issuers, at most $f$ are Byzantine and may choose arbitrary signed intervals. Up to $r$ further intervals from honest issuers may fail to cover the true state $z^\star$; honest measurement errors may be dependent. The fault count concerns trusted, enrolled issuing authorities, not merely distinct public keys. The gateway rejects repeated public keys and trust-domain labels, but this enforces the configured enrollment policy rather than proving the organizational independence of distinct registrants or defeating Sybil attacks. A missing or invalid signature from a required roster member causes the protocol to abort, rather than yielding a certified inconsistent decision.

Let $\ell_{ia}$ and $u_{ia}$ be the lower and upper endpoints of $U_i$ in coordinate $a$. The signed message digest is
\begin{equation}\label{eq:signedbox}
m_i=\operatorname{MiMC}(c,i,\ell_{i1},u_{i1},\ldots,\ell_{id},u_{id}).
\end{equation}
The circuit verifies an EdDSA signature on this digest. Its public statement includes $c$, $\Omega$, $P$, support threshold $k$, the decision code, and the ordered issuer public keys. The shared context binds the task, coordinate and calibration conventions, policy, threshold, issuer roster, report digest, challenge nonce, and expiry. The gateway reconstructs this statement from trusted configuration rather than accepting arbitrary task parameters from a vehicle.

Cryptographic authentication establishes which key signed which interval bounds, but not whether those bounds correctly enclose the physical event. Physical interpretation also requires sensor calibration, clock synchronization, coordinate registration, and defensible uncertainty envelopes. A continuous uncertainty region should be conservatively mapped to every quantization cell it intersects; the corresponding compliance inference requires those cells to lie wholly within the permitted region. The prototype assumes that this adaptation has been completed upstream, and rejects out-of-range coordinates instead of clamping them. A common systematic bias among honest sources can invalidate the coverage assumption without being detected by the proof.

The roster is complete and fixed for a task. If a required issuer withholds its evidence, certification may be blocked. Partial participation would require authentication of the active roster and adjusted fault and miscoverage allowances. Thus, the prototype targets task-level certification that tolerates some delay; it makes no availability guarantee for safety-critical real-time V2X control.
